\scanpage{8}\mspage{5 bis}

5. On peut aussi donner une condition d'\uline{existence} des $\pi_i$, liée aux conditions (4.4). Pour ceci, on suppose que $\unsure{E|\mathcal E}$ est contenu dans un sur-anneau gradué $F$ [p. ex. \inserted{l'anneau} des endom. d'un $M^*$ comme dans 1.] pour lequel les $\pi_i$ satisfont (1.1.) et les relations correspondantes à (1.3) (avec $E$ remplacé par $F$) sont satisfaites, \inserted{et (3.0) : $\pi_i=0$ si $i\notin[0,2n]$} Supposons que $\forall i$, la condition $(4.4)_i$ soit satisfaite, (avec $v_i,w_i\in E$ bien sur). On va en conclure que
\[
 \pi_i\in E\qquad\text{pour tt }i
 \tag*{(5.1)}
\]
\deleted{\unsure{Raisonnons prouvons en fait|Raisonnons Supposons provisoirement}} \inserted{supposons prouvé \deleted{\unsure{que}} déjà} $\pi_\alpha\in E$ si $\alpha\notin[i,2n-i]$, où $i<n$, et prouvons que l'on peut remplacer $i$ par $i+1$, \inserted{i.e $\pi_i$ et $\pi_{2n-i}\in E$.} Trivial si $i=n-1$, supposons donc $i\le n-2$.

Noter que, avec les notations du 3., \deleted{\unsure{on peut}} \inserted{on aura $\varphi_i,\psi_{2n-i}\in E$, \deleted{\unsure{d}} et} \deleted{\unsure{remplacer} $v_i$, \inserted{$w_i$} par} \inserted{donc} \deleted{\illegible{1}} $\psi_{2n-i}v_i\varphi_i,\ \varphi_iw_i\psi_{2n-i}\in E$, i.e
\[
 \left\{\begin{aligned}
 \inserted{v' =}\,\psi_{2n-i}v_i\varphi_i&=\pi_{2n-i}v_i\pi_i\in E\\
 w'_i=\varphi_i w_i\psi_{2n-i}&=\pi_i w_i\pi_{2n-i}\in E
 \end{aligned}\right.
\]
\deleted{et} \uline{co}mme \deleted{\unsure{on a}}
\[
 w'_iv'_i=\pi_i w_i\pi_{2n-i}v_i\pi_i
 \overset{\inserted{\text{car }w_i\pi_{2n-i}=\pi_i w_i}}{=}
 \pi_i\pi_i\unsure{W_i|w_i}v_i\pi_i
 \overset{\inserted{\text{car }w_iv_i\pi_i=\pi_i}}{=}\pi_i^3=\pi_i .
\]
et de $=$ \deleted{\unsure{$=\pi_i$}} $v'_i w'_i=\pi_{2n-i}$, on conclut que $\pi_i\in E$, $\pi_{2n-i}\in E$.
% archive sheet number: 7
